\chapter{Réalisation et résultats}
\markboth{Chapitre 4}{Réalisation et résultats}
\setcounter{chapter}{4}
\setcounter{section}{0}
\minitoc
\newpage

%% ═══════════════════════════════════════════════════════════════════════════════
\section{Introduction}

Ce chapitre présente l'environnement de développement, les détails
d'implémentation des composants clés, et les résultats quantitatifs obtenus
lors des collectes réelles effectuées au 2 mai 2026.

%% ═══════════════════════════════════════════════════════════════════════════════
\section{Environnement de développement}

\subsection{Stack technique}

\begin{table}[H]
  \centering
  \begin{tabular}{L{3cm} L{4cm} L{6cm}}
    \toprule
    \rowcolor{lightBlue!50}
    \textbf{Catégorie} & \textbf{Outil / Version} & \textbf{Rôle} \\
    \midrule
    Langage & Python 3.11 & Collecteurs, pipeline, API, NLP \\
    API Framework & FastAPI + Uvicorn & Exposition REST, BackgroundTasks \\
    Base de données & PostgreSQL 15-alpine & Stockage relationnel, tableaux natifs \\
    Pilote DB async & asyncpg 0.29 & Pool de connexions hautes performances \\
    HTTP client & httpx (async) & BOAMP, TED, Malt \\
    LinkedIn & linkedin-api 2.x & Scraping via mobile API interne \\
    Retry & tenacity & Backoff exponentiel sur APIs externes \\
    NLP & scikit-learn 1.4 & TF-IDF, LinearSVC, CalibratedCV \\
    LLM & Groq (llama-3.1-8b) & Génération messages/emails \\
    Automatisation & n8n 1.x & Workflows, scheduler, webhooks \\
    BI & Metabase OSS & Dashboard PostgreSQL \\
    Conteneurisation & Docker Compose v2 & 5 services en une commande \\
    Logging & Loguru & Logs structurés (INFO/SUCCESS/WARNING) \\
    OS & Windows 11 / Ubuntu & Développement / production \\
    \bottomrule
  \end{tabular}
  \caption{Environnement technique complet}
  \label{tab:env}
\end{table}

\subsection{Architecture Docker Compose}

La plateforme est entièrement conteneurisée. Un seul fichier
\texttt{docker-compose.yml} orchestre 5 services~:

\begin{figure}[H]
  \centering
  \begin{tikzpicture}[
    svc/.style={draw=#1!70, fill=#1!15, rounded corners=6pt,
                minimum width=3.2cm, minimum height=1.4cm,
                font=\small\bfseries\color{#1}, align=center, inner sep=8pt},
    port/.style={font=\scriptsize\color{espritGray}},
    dep/.style={-{Stealth[length=5pt]}, espritGray!50, dashed}
  ]
    \node[svc=espritBlue] (api) at (0, 0)
      {leadgen\_api\\{\small :8000}\\{\scriptsize Python 3.11}};
    \node[svc=darkGreen] (pg) at (-4, -2.5)
      {leadgen\_postgres\\{\small :5432}\\{\scriptsize PostgreSQL 15}};
    \node[svc=espritRed] (redis) at (-1.2, -2.5)
      {leadgen\_redis\\{\small :6379}\\{\scriptsize Redis 7}};
    \node[svc=darkBlue] (n8n) at (2, -2.5)
      {leadgen\_n8n\\{\small :5678}\\{\scriptsize n8n 1.x}};
    \node[svc=espritGray] (mb) at (5, -2.5)
      {leadgen\_metabase\\{\small :3000}\\{\scriptsize Metabase OSS}};

    \draw[dep] (api) -- (pg)
      node[midway, left, font=\scriptsize] {asyncpg};
    \draw[dep] (api) -- (redis)
      node[midway, font=\scriptsize] {};
    \draw[dep] (n8n) -- (pg);
    \draw[dep] (mb) -- (pg)
      node[midway, right, font=\scriptsize] {JDBC};
    \draw[dep] (n8n) |- ++(0, 0.8) -- (api)
      node[midway, above, font=\scriptsize] {HTTP REST};
  \end{tikzpicture}
  \caption{Architecture Docker Compose — 5 services}
  \label{fig:docker}
\end{figure}

%% ═══════════════════════════════════════════════════════════════════════════════
\section{Implémentation des collecteurs}

\subsection{Gestionnaire de sessions LinkedIn}

Pour éviter la détection par LinkedIn, un pool de 2 comptes est géré par
\texttt{collectors/session\_manager.py} avec rotation automatique.

\begin{lstlisting}[language=Python3, caption={Rotation de comptes LinkedIn (extrait)}, label=lst:session]
class SessionManager:
    async def initialize(self):
        for email, pwd in zip(self.emails, self.passwords):
            client = LinkedInClient(email, pwd)
            if await client.connect():
                self._clients.append(client)
                logger.success(f"Compte pret: {email}")
        logger.info(f"Pool: {len(self._clients)}/{len(self.emails)} actifs")

    def get_client(self) -> LinkedInClient:
        # Rotation round-robin
        client = self._clients[self._index % len(self._clients)]
        self._index += 1
        return client
\end{lstlisting}

\subsection{Déduplication dans le pipeline}

La requête de déduplication est au cœur de la fiabilité du système~:

\begin{lstlisting}[language=Python3, caption={Requête de déduplication 30 jours}, label=lst:dedup]
existing = await conn.fetchval("""
    SELECT id FROM opportunities
    WHERE company_id    = $1
      AND technologies  @> ARRAY[$2]
      AND opportunity_type = $3
      AND created_at    > NOW() - INTERVAL '30 days'
    LIMIT 1
""", company_id, prospect.technology, opportunity_type)

if existing:
    return False  # Doublon -- skip
\end{lstlisting}

\subsection{Collecteur TED EU — gestion du multilinguisme}

\begin{lstlisting}[language=Python3, caption={Extraction multilingue du nom acheteur TED}, label=lst:ted_impl]
async def _search(self, country_code: str, keyword: str) -> list[TedNotice]:
    query = f'notice-title="{keyword}" AND buyer-country={country_code}'
    payload = {"query": query, "fields": TED_FIELDS, "limit": 15}
    async with httpx.AsyncClient(timeout=30) as client:
        r = await client.post(TED_API_URL, json=payload)
        r.raise_for_status()
    notices = []
    for raw in r.json().get("notices", []):
        buyer = raw.get("buyer-name", {})
        # Preference : francais puis anglais
        name = buyer.get("fra") or buyer.get("eng") or "Acheteur EU"
        title = raw.get("notice-title", {})
        title_text = title.get("fra") or title.get("eng") or keyword
        notices.append(TedNotice(buyer_name=name, title=title_text, ...))
    return notices
\end{lstlisting}

%% ═══════════════════════════════════════════════════════════════════════════════
\section{API REST FastAPI}

\subsection{Endpoint de collecte asynchrone}

L'endpoint \texttt{POST /api/collect/all} retourne immédiatement pendant que
la collecte tourne en arrière-plan grâce à \texttt{BackgroundTasks} de FastAPI~:

\begin{lstlisting}[language=Python3, caption={Collecte déclenchée en BackgroundTask}, label=lst:collect_api]
@router.post("/collect/all")
async def trigger_full_collect(
    background_tasks: BackgroundTasks,
    _: str = Depends(verify_api_key)
):
    background_tasks.add_task(_run_all_collectors)
    return {"status": "started",
            "message": "Collecte complete demarree en arriere-plan"}

async def _run_all_collectors():
    from collectors.run_all import run_all
    await run_all(dry_run=False)
\end{lstlisting}

\subsection{Endpoint NLP}

\begin{lstlisting}[language=Python3, caption={Classification NLP via API}, label=lst:nlp_api]
@router.post("/classify")
async def classify_text(data: ClassifyRequest,
                        _: str = Depends(verify_api_key)):
    clf = get_classifier()
    result = clf.predict(data.text)
    return {
        "sector_label":         result.sector_label,
        "priority_label":       result.priority_label,
        "sector_confidence":    round(result.sector_confidence, 3),
        "priority_confidence":  round(result.priority_confidence, 3),
    }
\end{lstlisting}

%% ═══════════════════════════════════════════════════════════════════════════════
\section{Interfaces utilisateur}

\subsection{Dashboard Metabase}

Metabase se connecte directement à PostgreSQL (host~: \texttt{postgres},
port~: 5432). Les requêtes recommandées pour l'équipe commerciale~:

\begin{lstlisting}[language=SQL, caption={Top 20 prospects pour l'équipe commerciale}, label=lst:metabase]
SELECT c.name, o.opportunity_type, o.source_platform,
       o.country, o.priority_score, o.technologies[1] as tech,
       o.created_at
FROM opportunities o
JOIN companies c ON c.id = o.company_id
WHERE o.status = 'new'
ORDER BY o.priority_score DESC
LIMIT 20;
\end{lstlisting}

\subsection{Interface Swagger (FastAPI /docs)}

L'API expose une documentation Swagger interactive à l'adresse
\texttt{http://localhost:8000/docs}, permettant de tester tous les endpoints
directement depuis le navigateur.

%% ═══════════════════════════════════════════════════════════════════════════════
\section{Résultats et évaluation}

\subsection{Volume de données collectées}

Au 2 mai 2026, après les collectes initiales, la base de données contient
\textbf{815 prospects} répartis comme suit~:

\begin{figure}[H]
  \centering
  \begin{tikzpicture}
    \begin{axis}[
      ybar,
      width=13cm, height=6cm,
      bar width=22pt,
      ymin=0, ymax=260,
      ylabel={Nombre de prospects},
      xlabel={Source et pays},
      symbolic x coords={
        BOAMP-FR,
        TED-FR, TED-BE, TED-CH, TED-LU,
        LI-Jobs-FR, LI-Jobs-BE, LI-Jobs-TN, LI-Jobs-LU
      },
      xtick=data,
      x tick label style={font=\scriptsize, rotate=30, anchor=east},
      nodes near coords,
      nodes near coords style={font=\scriptsize\bfseries, rotate=90, anchor=west},
      enlarge x limits=0.08,
      legend style={at={(0.5,-0.35)}, anchor=north, font=\small, legend columns=2},
    ]
      \addplot+[fill=espritBlue!80] coordinates {
        (BOAMP-FR, 236)
      };
      \addplot+[fill=darkGreen!70] coordinates {
        (TED-FR, 145) (TED-BE, 105) (TED-CH, 98) (TED-LU, 74)
      };
      \addplot+[fill=espritRed!70] coordinates {
        (LI-Jobs-FR, 122) (LI-Jobs-BE, 12)
        (LI-Jobs-TN, 11) (LI-Jobs-LU, 10)
      };
      \legend{BOAMP (b2b\_tender), TED EU (b2b\_tender),
              LinkedIn Jobs (outsourcing\_signal)}
    \end{axis}
  \end{tikzpicture}
  \caption{Volume de prospects collectés par source et par pays (2 mai 2026)}
  \label{fig:results_volume}
\end{figure}

\begin{table}[H]
  \centering
  \begin{tabular}{L{2.8cm} C{1.5cm} C{1.8cm} C{1.8cm} C{3cm}}
    \toprule
    \rowcolor{lightBlue!50}
    \textbf{Source} & \textbf{Pays} & \textbf{Prospects} &
    \textbf{Score moy.} & \textbf{Type signal} \\
    \midrule
    BOAMP & FR & 236 & 66 & \texttt{b2b\_tender} \\
    TED EU & FR & 145 & 65 & \texttt{b2b\_tender} \\
    TED EU & BE & 105 & 63 & \texttt{b2b\_tender} \\
    TED EU & CH &  98 & 65 & \texttt{b2b\_tender} \\
    TED EU & LU &  74 & 66 & \texttt{b2b\_tender} \\
    LinkedIn Jobs & FR & 122 & 79 & \texttt{outsourcing\_signal} \\
    LinkedIn Jobs & BE &  12 & 83 & \texttt{outsourcing\_signal} \\
    LinkedIn Jobs & TN &  11 & 75 & \texttt{outsourcing\_signal} \\
    LinkedIn Jobs & LU &  10 & 84 & \texttt{outsourcing\_signal} \\
    LinkedIn Jobs & CH &   2 & 87 & \texttt{outsourcing\_signal} \\
    \midrule
    \textbf{TOTAL} & & \textbf{815} & \textbf{$\approx$70} & \\
    \bottomrule
  \end{tabular}
  \caption{Résultats détaillés par source (au 2 mai 2026)}
  \label{tab:results}
\end{table}

\subsection{Répartition par niveau de priorité}

\begin{figure}[H]
  \centering
  \begin{tikzpicture}
    \begin{axis}[
      xbar,
      width=11cm, height=4cm,
      bar width=20pt,
      xmin=0, xmax=450,
      xlabel={Nombre de prospects},
      symbolic y coords={Score 0--59, Score 60--79, Score 80--100},
      ytick=data,
      y tick label style={font=\small},
      nodes near coords,
      enlarge y limits=0.4,
    ]
      \addplot+[fill=espritGray!50] coordinates {(270,{Score 0--59})};
      \addplot+[fill=espritBlue!70] coordinates {(387,{Score 60--79})};
      \addplot+[fill=darkGreen!80]  coordinates {(158,{Score 80--100})};
    \end{axis}
  \end{tikzpicture}
  \caption{Distribution des prospects par plage de score}
  \label{fig:score_dist}
\end{figure}

\subsection{Scores moyens par pays (LinkedIn Jobs)}

Les prospects LinkedIn Jobs affichent des scores significativement plus élevés
car ils cumulent souvent pays premium (LU/CH) et technologies rares (COBOL).

\begin{figure}[H]
  \centering
  \begin{tikzpicture}
    \begin{axis}[
      ybar,
      width=8cm, height=5cm,
      bar width=20pt,
      ymin=60, ymax=95,
      ylabel={Score moyen},
      symbolic x coords={FR, BE, TN, LU, CH},
      xtick=data,
      nodes near coords,
      nodes near coords style={font=\small\bfseries},
      enlarge x limits=0.25,
    ]
      \addplot+[fill=espritRed!70] coordinates {
        (FR, 79) (BE, 83) (TN, 75) (LU, 84) (CH, 87)
      };
    \end{axis}
  \end{tikzpicture}
  \caption{Score moyen par pays — collecteur LinkedIn Jobs}
  \label{fig:avg_score_country}
\end{figure}

\subsection{Performances du module NLP}

\begin{table}[H]
  \centering
  \begin{tabular}{L{5cm} C{2.5cm} C{2.5cm}}
    \toprule
    \rowcolor{lightBlue!50}
    \textbf{Métrique} & \textbf{Classifieur secteur} & \textbf{Classifieur priorité} \\
    \midrule
    Cross-validation 5-fold & \textbf{75.0\%} & -- \\
    Accuracy (test set) & 70.2\% & 85.0\% \\
    F1-macro & 68\% & 81\% \\
    Latence prédiction & $\approx$100\,ms & $\approx$100\,ms \\
    Taille modèle (.joblib) & $\approx$8\,MB & $\approx$6\,MB \\
    \bottomrule
  \end{tabular}
  \caption{Performances des modèles NLP}
  \label{tab:nlp_perf}
\end{table}

\subsection{Performances du pipeline end-to-end}

\begin{figure}[H]
  \centering
  \begin{tikzpicture}[
    node distance=0.4cm,
    pipstep/.style={draw=darkBlue!60, fill=pipColor, rounded corners=4pt,
                 minimum width=8cm, minimum height=0.7cm,
                 font=\small, align=center, inner sep=6pt},
    time/.style={draw=espritRed!50, fill=lightAmber!30, rounded corners=3pt,
                 minimum width=2.8cm, minimum height=0.7cm,
                 font=\scriptsize\bfseries\color{espritRed}, align=center},
    arr/.style={-{Stealth[length=5pt]}, darkBlue!60, thick}
  ]
    \node[pipstep] (s0) at (0,  0.00) {Source externe (API/scraping)};
    \node[time, right=0.2cm of s0] {0.3s--2.0s};
    \node[pipstep] (s1) at (0, -0.95) {Collecteur Python (normalisation)};
    \node[time, right=0.2cm of s1] {$<$5\,ms};
    \node[pipstep] (s2) at (0, -1.90) {pipeline/ingest.py (upsert + dédup + scoring)};
    \node[time, right=0.2cm of s2] {\textasciitilde 50\,ms};
    \node[pipstep] (s3) at (0, -2.85) {PostgreSQL INSERT + indexation};
    \node[time, right=0.2cm of s3] {\textasciitilde 5\,ms};
    \node[pipstep] (s4) at (0, -3.80) {/api/classify (NLP optionnel)};
    \node[time, right=0.2cm of s4] {\textasciitilde 100\,ms};
    \node[pipstep] (s5) at (0, -4.75) {/api/generate/message (Groq LLM)};
    \node[time, right=0.2cm of s5] {1--3\,s};
    \draw[arr] (s0.south) -- (s1.north);
    \draw[arr] (s1.south) -- (s2.north);
    \draw[arr] (s2.south) -- (s3.north);
    \draw[arr] (s3.south) -- (s4.north);
    \draw[arr] (s4.south) -- (s5.north);
  \end{tikzpicture}
  \caption{Latences du pipeline de traitement bout-en-bout}
  \label{fig:latencies}
\end{figure}

\subsection{Exemple de prospect généré}

\begin{successbox}[Exemple de sortie --- BNP Paribas COBOL France]
\textbf{Prospect ingéré~:}
\begin{itemize}[nosep]
  \item Entreprise~: BNP Paribas | Technologie~: COBOL | Pays~: France
  \item Type~: \texttt{outsourcing\_signal} | Score~: \textbf{95/100}
  \item Secteur NLP~: Finance/Banque (confiance 0.92) | Priorité~: HIGH (0.88)
\end{itemize}
\vspace{0.3cm}
\textbf{Message LinkedIn généré par Groq~:}
\begin{quote}
\itshape
Bonjour [Prénom],

J'ai remarqué que BNP Paribas recrute activement des profils COBOL en ce moment.

Chez Solvinya Group, nous disposons d'une équipe COBOL/Mainframe disponible
immédiatement — nous intervenons souvent quand le recrutement prend trop de
temps ou que la mission est trop courte pour justifier un CDI.

Vous travaillez sur une migration ou une maintenance de système legacy~?
\end{quote}
\end{successbox}

%% ═══════════════════════════════════════════════════════════════════════════════
\section{Tests}

\begin{table}[H]
  \centering
  \begin{tabular}{L{3.5cm} L{4cm} L{5cm}}
    \toprule
    \rowcolor{lightBlue!50}
    \textbf{Module testé} & \textbf{Type de test} & \textbf{Résultat} \\
    \midrule
    API REST & pytest-asyncio (endpoints) & Health, ingest, classify, generate $\checkmark$ \\
    NLP Classifier & Accuracy CV-5 & 75.0\% $\geq$ 75\% seuil $\checkmark$ \\
    Pipeline ingest & Déduplication & Doublon correctement ignoré $\checkmark$ \\
    Scoring & Calculs unitaires & COBOL/LU = 100 $\checkmark$ \\
    TED collector & Dry-run & 422 prospects validés $\checkmark$ \\
    BOAMP collector & Dry-run & 236 prospects validés $\checkmark$ \\
    LinkedIn B2B & Dry-run & Crédit Agricole, UBCI détectés $\checkmark$ \\
    \bottomrule
  \end{tabular}
  \caption{Tableau de synthèse des tests}
  \label{tab:tests}
\end{table}

%% ═══════════════════════════════════════════════════════════════════════════════
\section{Conclusion}

Ce chapitre a démontré la pleine fonctionnalité de la plateforme LeadGen 360+~:
815 prospects collectés sur 3 sources actives, modèle NLP atteignant le seuil de
75\%, messages de prospection générés en 1 à 3 secondes. Le système est
conteneurisé, documenté et automatisé, prêt pour une utilisation en production
par l'équipe commerciale de Solvinya Group.
